% DELTA23: exact fixed-point coordinates, orthogonal group, and outer signs.
\subsection{The orthogonal geometry of the new Hurwitz block}

\paragraph{D23.C23 (an explicit minus-type quadratic geometry).}
Take block zero of the degree-$1344$ Hurwitz action, with its
$112$ labels in the order recorded in
\path{agents/five_level_six/hurwitz_block_stabilizer.npz}.
For a tuple entry $g_i=(v_i,A_i)\in
\mathbb F_3^2\rtimes\operatorname{SL}_2(\mathbb F_3)$, define
its unique fixed point by
\[
                     q_i=(I-A_i)^{-1}v_i.
\]
All these inverses exist. Normalize the ordered linear tuple by
an element $B\in\operatorname{SL}_2(\mathbb F_3)$ to
\[
 A_1=A_2=\begin{pmatrix}0&1\\2&1\end{pmatrix},\qquad
 A_3=A_4=\begin{pmatrix}1&1\\2&0\end{pmatrix}.
\tag{D23.C23.1}
\]
There are exactly two such normalizers, differing by $-I$.
Thus
\[
 [w]=[\,B(q_1-q_4),B(q_2-q_4),B(q_3-q_4)\,]
                  \in\mathbb P^5(\mathbb F_3)
\tag{D23.C23.2}
\]
is well defined. These $112$ projective vectors are distinct
and are exactly the singular points of the form $w^tQw$, where
\[
 Q=\begin{pmatrix}
 0&1&1&1&1&0\\
 1&1&0&1&0&1\\
 1&0&1&1&0&2\\
 1&1&1&0&1&1\\
 1&0&0&1&1&2\\
 0&1&2&1&2&0
 \end{pmatrix}.
\tag{D23.C23.3}
\]
The form has rank six and determinant one. Hence it has minus
type over $\mathbb F_3$ and Witt index two.

Every block-stabilizer permutation has a unique projective
linear realization on these coordinates. In particular each
of the thirteen recorded stabilizer generators lifts to a
matrix $M$ satisfying
\[
                         M^tQM=Q,\qquad\det M=1.
\tag{D23.C23.4}
\]
The complete coordinate table and all thirteen matrices are
in \path{agents/characters/records/hurwitz_orthogonal112.npz}
and the corresponding JSON record. Consequently the independently
computed block image of order $6531840$ is exactly
\[
 H=\operatorname{PSO}_6^-(3),\qquad
 H'=\operatorname{P\Omega}_6^-(3)\simeq
                  \operatorname{PSU}_4(3),
 \qquad |H'|=3265920.
\tag{D23.C23.5}
\]
In particular $H'$ is nonabelian simple. The space of
$H$-invariant symmetric bilinear forms on this six-dimensional
space has dimension one. Thus any nonzero symmetric form
obtained independently from a cohomological construction and
preserved by these same generators is a scalar multiple of $Q$.
This last assertion identifies multipliers without requiring
an arbitrary scalar normalization of that second form.

\begin{proof}
The affine conjugation formula sends $q_i$ to $u+Bq_i$;
subtracting $q_4$ removes $u$. The simultaneous centralizer of
the ordered tuple (D23.C23.1) in $\operatorname{SL}_2(3)$ is
$\{I,-I\}$, verified directly from its two distinct matrices.
Searching the twenty-four matrices of $\operatorname{SL}_2(3)$
therefore gives the stated projective normalization, with no
choice left except the harmless sign. The certificate checks
the existence of precisely these two transporters for every
one of the $112$ tuples.

The program \path{agents/characters/code/hurwitz_orthogonal112.py}
performs only exact arithmetic over $\mathbb F_3$.
The $112$ equations $w^tQw=0$ in the twenty-one coefficients
of a quadratic form have a one-dimensional solution space;
its normalized matrix is (D23.C23.3). Exact row reduction gives
rank six and determinant one. It then enumerates all
$(3^6-1)/2=364$ projective points and checks that its singular
locus is exactly the recorded set of $112$ points. For an
even-dimensional form of dimension $2m$, the split case has
$(-1)^m\det Q$ square. Here $m=3$ and $-\det Q=2$ is a
nonsquare, giving minus type.

For each stabilizer permutation choose six independent
coordinate vectors. Their images determine a linear map up
to six individual signs. Fixing one sign leaves only
$2^5=32$ possibilities. Exactly one of these maps induces the
given permutation on all $112$ points. It is invertible and
satisfies (D23.C23.4). The same procedure for the identity
permutation confirms faithfulness modulo scalar matrices.
Finally the simultaneous equations $M^tQ'M=Q'$ on the
twenty-one symmetric coefficients have dimension one, proving
the invariant-form assertion independently of the singular
locus computation.

The standard orthogonal order formula gives
\[
 |O_6^-(3)|=2\cdot3^6(3^3+1)(3^2-1)(3^4-1)=26127360.
\]
Taking determinant one and then dividing by $\{I,-I\}$
gives $|\operatorname{PSO}_6^-(3)|=6531840$.
The computed inclusion and equality of orders prove the
first group identification. The spinor computation in C24
is onto on this group and trivial on its derived subgroup;
its projective kernel has order $3265920$, the independently
computed order of $H'$. Hence $H'=\operatorname{P\Omega}_6^-(3)$.
The exceptional isomorphism with $\operatorname{PSU}_4(3)$
and simplicity are the standard finite classical-group facts
for this exact parameter, not an inference from order alone.
For explicit primary reference conventions, the GAP Character
Table Library entry \texttt{U4(3)} lists both
\texttt{POmega(-1,6,3)} and \texttt{PSU(4,3)} as its group
constructions; the ATLAS entry gives its order and the
$112$-point action.
\footnote{Primary pages checked in this run:
\url{https://www.math.rwth-aachen.de/~Thomas.Breuer/ctbllib/ctbltoc/data/U4(3).html}
and \url{https://brauer.maths.qmul.ac.uk/Atlas/clas/U43/}.
Cached pages and exact locators are in
\path{agents/characters/records/orthogonal_primary_sources.txt}.}
\end{proof}

\paragraph{D23.C24 (the full normalizer and its two visible outer signs).}
Write $S=H'$. In this statement
\[
 GO^*(Q)=\{M\in GL_6(3):M^tQM=\mu(M)Q,
                                 \ \mu(M)\in\mathbb F_3^\times\}
\]
includes nonsquare similitudes, and $PGO^*(Q)$ is its quotient
by the two scalar matrices. Then
\[
 N_{\mathfrak S_{112}}(H)=N_{\mathfrak S_{112}}(S)
       =PGO^*(Q),\qquad |PGO^*(Q)|=26127360,
\tag{D23.C24.1}
\]
and the maps determinant and multiplier give
\[
 PGO^*(Q)/H\simeq C_2\times C_2.
\tag{D23.C24.2}
\]
Both are well defined projectively, since the dimension is
even and both scalar squares are one. Furthermore
\[
                   PGO^*(Q)/S\simeq D_8,
\]
and its central subgroup of order two is exactly $H/S$.

The spinor norms of the thirteen stabilizer matrices, in
their recorded order, are
\[
                       1,1,1,2,2,2,2,2,2,2,2,2,2.
\tag{D23.C24.3}
\]
Here $2=-1$ is the nonsquare class. The norm of $-I$ is one,
so this character descends to $H$ and has kernel $S$.
Explicit matrices and permutations for the two other outer
generators are in
\path{agents/characters/records/hurwitz_orthogonal_signs.npz}.
One is a reflection with $(\det,\mu)=(2,1)$; the other has
$(\det,\mu)=(1,2)$. Their commutator has determinant and
multiplier one but spinor norm two.

\begin{proof}
The exact point-stabilizer orbit sizes are $1,30,81$ for
both $H$ and $S$. The orbit of length thirty is precisely
orthogonality for (D23.C23.3). A permutation normalizing either
group must preserve that uniquely distinguished orbital graph.
The full automorphism group of this graph is $PGO^*(Q)$;
this is the standard $O_6^-(3)$ graph with parameters
$(112,30,2,10)$. Brouwer's primary graph record identifies
its full automorphism group, including the nonsquare
similitudes, as $U_4(3).D_8$ of order $26127360$.
\footnote{A. E. Brouwer, \emph{The $O_6^-(3)$ graph},
\url{https://aeb.win.tue.nl/graphs/U4_3.html}, Group and
Construction paragraphs, checked in this run. Its star
notation is retained here to avoid confusing isometries with
all similitudes.}
Conversely every similitude normalizes $SO(Q)$ and its
spinor kernel, proving both equalities in (D23.C24.1).

For the exact spinor certificate, reflect in a nonsingular
vector $v$ by
\[
 R_v=I-\frac{2}{v^tQv}\,v(v^tQ).
\]
An orthogonal Gram--Schmidt basis reduces each isometry to
the identity by these reflections. If $Mv-v$ is nonsingular,
reflect in it; otherwise reflect successively in $Mv+v$
and $v$. The latter two vectors are nonsingular because
$v$ is and $Mv-v$ is singular. Multiplying their quadratic
norms gives the spinor norm. The exact replay
\path{agents/characters/code/hurwitz_orthogonal_signs.py}
checks every reflection and the resulting identity matrix,
the determinant parity, (D23.C24.3), and the norm of $-I$.

The same certificate constructs a reflection $R$ and a
nonsquare similitude $M$ with the determinant and multiplier
values stated above. Their squares belong to $S$, while
their commutator belongs to $H\setminus S$.
These two independent outer cosets give (D23.C24.2).
For any similitude $C$ of multiplier $\mu$,
$CR_vC^{-1}=R_{Cv}$ and
$Q(Cv)=\mu Q(v)$. An element of $SO(Q)$ is a product of an
even number of reflections; hence conjugation by $C$ preserves
its spinor norm. It follows that $H/S$ is central in the
outer quotient. The nontrivial commutator of $R,M$, together
with their involutory images, identifies this order-eight
quotient as $D_8$.
\end{proof}

\paragraph{D23.C24a (one common outer type across the twelve blocks).}
Suppose the independently computed full Hurwitz group has
socle $S^{12}$, with its twelve factors supported on the
twelve blocks, and its local transition maps lie in $H$
after the chosen group-transversal identifications.
Any permutation normalizing this full group has local maps
in $PGO^*(Q)$, and all twelve local maps have the same
determinant/multiplier pair in $PGO^*(Q)/H$.
Indeed it permutes the simple factors and their blocks.
Conjugating a group transition from block $i$ to block $j$
shows that the two local maps have equal cosets modulo $H$.
The group is transitive on the blocks, so these cosets are
constant. This is a necessary condition on the actual GT
conjugator supplied by the Hurwitz functoriality theorem.
The identification of this common pair with the original
harmonic and cubic frame bits is a further calculation;
it is not asserted by the orthogonal group identification alone.
